\documentclass[11pt]{article}
\usepackage{mathblatt}
\usepackage{adjustbox,varwidth}
% gesetzt von werkzeuge/setzer.py; Lösungen nach bau/bauregeln.md „Lösungen“.
\newcommand{\kennung}{T74}
\begin{document}
\pfheftstil
\fancyfoot[C]{{\footnotesize\color{mbgrau}\kennung\hfill\thepage}}
\pfheftkopf{Säulen-, Balken- und Liniendiagramme}{Klasse 6 \textperiodcentered{} Lösungen}

\begin{pfloesung}{}
\lz{1b)}{\mbox{a) $10$} \mbox{b) $25$} \mbox{c) $0{,}1$}}{a) $50 : 5 = 10$; b) $100 : 4 = 25$; c) $1 : 10 = 0{,}1$}{}
\end{pfloesung}
\begin{pfloesung}{}
\lz{2.}{Ben $6$, Cem $14$ Bücher}{ein Kästchen: $10 : 5 = 2$; Ben: $3 \cdot 2 = 6$; Cem: $10 + 2 \cdot 2 = 14$}{}
\end{pfloesung}
\begin{pfloesung}{}
\lz{3.}{\mbox{a) $35\,000$} \mbox{b) $120\,000$} \mbox{c) $2\,500$}}{a) $35 \cdot 1000 = 35\,000$; b) $120 \cdot 1000 = 120\,000$; c) $2{,}5 \cdot 1000 = 2\,500$}{}
\end{pfloesung}
\begin{pfloesung}{}
\lz{4.}{Ja, im Juli $13\,000$ Gäste}{Ja; im Juli kamen $13 \cdot 1000 = 13\,000$ Gäste, mehr als $12\,000$.}{}
\end{pfloesung}
\begin{pfloesung}{}
\lz{5b)}{$50$ Punkte}{Gelb $90$, Blau $40$: $90 - 40 = 50$}{}
\end{pfloesung}
\begin{pfloesung}{}
\lz{6.}{Montag und Freitag}{Mo $175$, Fr $200$; Dienstag mit genau $150$ zählt nicht.}{}
\end{pfloesung}
\begin{pfloesung}{}
\lz{7.}{$150$\,€ (Mi, Fr, Sa)}{$150$\,€; mehr als $120$ Karten: Mi $140$, Fr $180$, Sa $200$ (Di mit genau $120$ zählt nicht); $3 \cdot 50 = 150$}{}
\end{pfloesung}
\begin{pfloesung}{}
\lz{8b)}{$6$\,°C; $7$\,°C wärmer}{Mittwoch $6$\,°C; $13 - 6 = 7$\,°C wärmer}{}
\end{pfloesung}
\begin{pfloesung}{}
\lz{9.}{2022, $10$ weniger}{2022: $30 - 20 = 10$ Mitglieder weniger}{}
\end{pfloesung}
\begin{pfloesung}{}
\lz{10.}{im Juni (um $30$\,cm gesunken)}{Im Juni: $70 - 40 = 30$\,cm, mehr als $20$\,cm. Im Mai sank er nur um $90 - 70 = 20$\,cm.}{}
\end{pfloesung}
\begin{pfloesung}{}
\lz{11b)}{Balken $9$ Kästchen lang (bis $45$)}{Balken Mi $9$ Kästchen lang; ein Kästchen: $20 : 4 = 5$; $45 : 5 = 9$\par\smallskip \adjustbox{max width=\linewidth,max height=4cm}{\begin{varwidth}{50cm}\begin{tikzpicture}[x=0.0700cm,y=0.55cm,line width=0.6pt]\draw[mbgitter,line width=0.3pt] (0,0) -- (0,4.4);\draw[mbgitter,line width=0.3pt] (5,0) -- (5,4.4);\draw[mbgitter,line width=0.3pt] (10,0) -- (10,4.4);\draw[mbgitter,line width=0.3pt] (15,0) -- (15,4.4);\draw[mbgitter,line width=0.3pt] (20,0) -- (20,4.4);\draw[mbgitter,line width=0.3pt] (25,0) -- (25,4.4);\draw[mbgitter,line width=0.3pt] (30,0) -- (30,4.4);\draw[mbgitter,line width=0.3pt] (35,0) -- (35,4.4);\draw[mbgitter,line width=0.3pt] (40,0) -- (40,4.4);\draw[mbgitter,line width=0.3pt] (45,0) -- (45,4.4);\draw[mbgitter,line width=0.3pt] (50,0) -- (50,4.4);\draw[mbgitter,line width=0.3pt] (55,0) -- (55,4.4);\draw[mbgitter,line width=0.3pt] (60,0) -- (60,4.4);\draw[black!45,line width=0.3pt] (0,0) -- (0,4.4);\node[below,font=\small,inner sep=3pt] at (0,0) {0};\draw[black!45,line width=0.3pt] (20,0) -- (20,4.4);\node[below,font=\small,inner sep=3pt] at (20,0) {20};\draw[black!45,line width=0.3pt] (40,0) -- (40,4.4);\node[below,font=\small,inner sep=3pt] at (40,0) {40};\draw[black!45,line width=0.3pt] (60,0) -- (60,4.4);\node[below,font=\small,inner sep=3pt] at (60,0) {60};\filldraw[fill=mbkasten,draw=black,line width=0.7pt] (0,3.70) rectangle (30,4.30);\node[left,font=\small,inner sep=4pt] at (0,4) {Mo};\filldraw[fill=mbkasten,draw=black,line width=0.7pt] (0,2.70) rectangle (50,3.30);\node[left,font=\small,inner sep=4pt] at (0,3) {Di};\filldraw[fill=mbkasten,draw=black,line width=0.7pt] (0,1.70) rectangle (45,2.30);\node[left,font=\small,inner sep=4pt] at (0,2) {Mi};\filldraw[fill=mbkasten,draw=black,line width=0.7pt] (0,0.70) rectangle (25,1.30);\node[left,font=\small,inner sep=4pt] at (0,1) {Do};\draw[-{Stealth[length=2mm]}] (0,0) -- (72.000,0) node[below right,font=\small,inner sep=2pt] {Gäste};\draw[-{Stealth[length=2mm]}] (0,0) -- (0,4.7);\end{tikzpicture}\end{varwidth}}}{}
\end{pfloesung}
\begin{pfloesung}{}
\lz{12.}{B $105$, C $75$ Stimmen}{ein Kästchen: $60 : 4 = 15$; B: $7 \cdot 15 = 105$; C: $5 \cdot 15 = 75$}{}
\end{pfloesung}
\begin{pfloesung}{}
\lz{13.}{\mbox{a) Schwimmen} \mbox{b) $3$ Kästchen hoch}}{a) Schwimmen: ein Kästchen $240 : 8 = 30$; $6 \cdot 30 = 180$; b) $90 : 30 = 3$ Kästchen\par\smallskip \adjustbox{max width=\linewidth,max height=4cm}{\begin{varwidth}{50cm}\saeulenab[ymax=9,ystep=1,yfein=1,ylabel=Mitglieder,hoehe=1.9,breite=0.8,ohnezahlen]{Fußball/8,Schwimmen/6,Tennis/3}\end{varwidth}}}{}
\end{pfloesung}
\begin{pfloesung}{}
\lz{14.}{$2\,500$ Fahrgäste}{$2\,500$; ein Kästchen: $1 : 2 = 0{,}5$; $2 + 0{,}5 = 2{,}5$; $2{,}5 \cdot 1000 = 2\,500$}{}
\end{pfloesung}
\begin{pfloesung}{}
\lz{15.}{Di und Do; $40$ mehr}{Dienstag und Donnerstag ($90$ und $70$; Montag mit genau $60$ zählt nicht); $90 - 50 = 40$ Brötchen mehr}{}
\end{pfloesung}
\begin{pfloesung}{}
\lz{16.}{im März, $3$\,kg}{im März: $7 - 4 = 3$\,kg}{}
\end{pfloesung}
\begin{pfloesung}{}
\lz{17.}{\mbox{a) $8$ Kästchen hoch} \mbox{b) $45$\,€ mehr}}{a) ein Kästchen: $90 : 6 = 15$; $120 : 15 = 8$ Kästchen; b) Di: $5 \cdot 15 = 75$; $120 - 75 = 45$\,€\par\smallskip \adjustbox{max width=\linewidth,max height=4cm}{\begin{varwidth}{50cm}\saeulenab[ymax=9,ystep=1,yfein=1,ylabel=Einnahmen,hoehe=1.9,breite=0.8,ohnezahlen]{Mo/6,Di/5,Sa/8}\end{varwidth}}}{}
\end{pfloesung}
\end{document}
